\section{Aprender la recompensa a partir de preferencias humanas}

\subsection{Comparación de preferencias}

Pedirle directamente a un ser humano una calificación absoluta es difícil y produce mucho ruido. Una mejor manera: mostrar dos opciones y dejar que el ser humano elija cuál es mejor.

$$p(y_1 \succ y_2 \mid x) = \sigma\left(R_\phi(x, y_1) - R_\phi(x, y_2)\right)$$

\begin{importantbox}{Modelo de Bradley-Terry}
Dados los datos de preferencia humana $(y_w, y_l)$ ($y_w$ es el preferido), se entrena el modelo de recompensa $R_\phi$ para maximizar:
$$\mathcal{L}(\phi) = \mathbb{E}\left[\log \sigma\left(R_\phi(x, y_w) - R_\phi(x, y_l)\right)\right]$$

Esta es la fórmula central del \textbf{entrenamiento del modelo de recompensa} en RLHF.
\end{importantbox}

\subsection{Pipeline de RLHF}

\begin{enumerate}
    \item Recopilar datos de comparación de preferencias humanas
    \item Entrenar el modelo de recompensa $R_\phi$
    \item Optimizar la política con RL (por ejemplo, PPO) para maximizar la salida de $R_\phi$
\end{enumerate}

\begin{warningbox}{El modelo de recompensa también puede ser explotado}
De manera similar al goal classifier, si la optimización de RL es excesiva, la política puede encontrar vulnerabilidades del modelo de recompensa en lugar de satisfacer realmente las preferencias humanas. Por eso normalmente se agrega una restricción de KL que limita qué tanto puede alejarse la política de la política inicial.
\end{warningbox}

\subsection{Resumen de la sección}

Aprender un modelo de recompensa a partir de comparaciones de preferencias es la base de RLHF. El modelo de Bradley-Terry convierte la comparación de preferencias en un sigmoide de la diferencia de recompensas, y su entrenamiento es simple y eficiente.

